%\ifodd\value{page}\hbox{}\newpage\fi
\chapter*{Dhyanam 4}

\begin{sloka}
{\sanskritfont
\begin{tabular}{c}
सकुङ्कुमविलेपनामलिकचुम्बिकस्तूरिकां\\
समन्दहसितेक्षणां सशरचापपाशाङ्कुशाम् ।\\
अशेषजनमोहिनीं अरुणमाल्यभूषाम्बरां\\
जपाकुसुमभासुरां जपविधौ स्मरेदम्बिकाम् ॥
\end{tabular}

\vspace{1.5em}

{\small \iastfont
sakuṅkumavilepanāmalikacumbikastūrikāṃ\\
samandahasitekṣaṇāṃ saśaracāpapāśāṅkuśām |\\
aśeṣajanamohinīṃ aruṇamālyabhūṣāmbarāṃ\\
japākusumabhāsurāṃ japavidhau smared ambikām ||\\
}

\vspace{1em}

{\small
\anvaya{%
जपविधौ सकुङ्कुमविलेपनाम् अलिकचुम्बिकस्तूरिकाम्
समन्दहसितेक्षणाम् सशरचापपाशाङ्कुशाम्
अशेषजनमोहिनीम् अरुणमाल्यभूषाम्बराम्
जपाकुसुमभासुराम् अम्बिकाम् स्मरेत् ॥
}
}

\small \iastfont \itmean{%
«Nel rito del japa si ricordi la Madre:
cosparsa di vermiglio, con il muschio che sfiora la fronte;
dallo sguardo dolcemente sorridente;
che regge frecce, arco, laccio e pungolo;
che incanta tutti gli esseri;
ornata e vestita di rosso;
splendente come il fiore di japa.»%
}

}
\end{sloka}

\wordmeaning{%
\wordline{सकुङ्कुम-विलेपनाम्}{sakuṅkuma-vilepanām}
  { cosparsa di vermiglio;}%
\wordline{अलिक-चुम्बि-कस्तूरिकाम्}{alikacumbi-kastūrikām}
  { con il muschio che sfiora la fronte;}%
\wordline{समन्द-हसित-ईक्षणाम्}{samanda-hasitekṣaṇām}
  { dallo sguardo dolcemente sorridente;}%
\wordline{स-शर-चाप-पाश-अङ्कुशाम्}{saśara-cāpa-pāśāṅkuśām}
  { che regge frecce, arco, laccio e pungolo;}%
\wordline{अशेष-जन-मोहिनीम्}{aśeṣa-jana-mohinīm}
  { che incanta tutti gli esseri;}%
\wordline{अरुण-माल्य-भूषा-अम्बराम्}{aruṇa-mālya-bhūṣāmbarām}
  { ornata e vestita di rosso;}%
\wordline{जपाकुसुम-भासुराम्}{japākusuma-bhāsurām}
  { splendente come il fiore di japa;}%
\wordline{जप-विधौ}{japa-vidhau}
  { nel rito del japa;}%
\wordline{स्मरेत्}{smaret}
  { si ricordi, si mediti;}%
\wordline{अम्बिकाम्}{ambikām}
  { la Madre divina;}%
}
